%! Special Directions --- Setting Up Eigenvalues — Unit 6, Lesson 6.0 — Homework
\documentclass[10pt]{article}
\usepackage{linalg-article}
\usepackage{linalg-boxes}
\newcommand{\TallMath}[1]{$\displaystyle #1\rule[-1.4em]{0pt}{3.2em}$}
\newcommand{\xx}{\mathbf{x}}
\newcommand{\yy}{\mathbf{y}}
\newcommand{\vv}{\mathbf{v}}
\newcommand{\zero}{\mathbf{0}}

\begin{document}

\pageheader{Unit 6, Lesson 6.0}{Homework}
\namedateperiod

\vspace{0.12in}

\begin{notesbox}{Practice}
The test is always the same: compute $A\xx$ and ask whether $A\xx=\lambda\xx$ for some number
$\lambda$. If yes, $\xx$ is an \textbf{eigenvector} and $\lambda$ is its \textbf{eigenvalue}.

\begin{enumerate}[leftmargin=1.9em, itemsep=12pt]
  \item \textbf{Test each candidate.} For $A=\begin{bmatrix} 2 & 1 \\ 1 & 2 \end{bmatrix}$, compute
        $A\xx$ and state whether $\xx$ is an eigenvector; if so, give $\lambda$.
    \begin{enumerate}[label=(\alph*), leftmargin=1.8em, itemsep=4pt]
      \item $\xx=\begin{bmatrix} 1 \\ 1 \end{bmatrix}$ \hfill
      \item $\xx=\begin{bmatrix} 1 \\ -1 \end{bmatrix}$ \hfill
      \item $\xx=\begin{bmatrix} 2 \\ 0 \end{bmatrix}$
    \end{enumerate}
    \vspace{1.2cm}
  \item \textbf{Find $\lambda$ from a given eigenvector.} Each $\xx$ below \emph{is} an eigenvector of
        the matrix beside it. Compute $A\xx$ and read off $\lambda$.
    \begin{enumerate}[label=(\alph*), leftmargin=1.8em, itemsep=4pt]
      \item $A=\begin{bmatrix} 5 & 0 \\ 0 & 3 \end{bmatrix}$, $\xx=\begin{bmatrix} 0 \\ 1 \end{bmatrix}$ \hfill
      \item $A=\begin{bmatrix} 6 & 2 \\ 0 & 4 \end{bmatrix}$, $\xx=\begin{bmatrix} 1 \\ -1 \end{bmatrix}$
    \end{enumerate}
    \vspace{1.1cm}
  \item \textbf{A flip: negative $\lambda$.} For $A=\begin{bmatrix} 1 & 2 \\ 2 & 1 \end{bmatrix}$,
        compute $A\begin{bmatrix} 1 \\ -1 \end{bmatrix}$. Find $\lambda$. What does the \emph{sign} of
        $\lambda$ say happened to the direction? \writelines{2}
  \item \textbf{The $\lambda=0$ direction (§5).} For $A=\begin{bmatrix} 3 & 6 \\ 1 & 2 \end{bmatrix}$,
        compute $A\begin{bmatrix} 2 \\ -1 \end{bmatrix}$. What is $\lambda$? Compute $\det A$. Is $A$
        invertible? Connect your two answers. \writelines{2}
  \item \textbf{Explain.} What does it mean for $\xx$ to be an eigenvector of $A$ with eigenvalue
        $\lambda$? Why is the \emph{zero} vector never counted as an eigenvector? \writelines{2}
\end{enumerate}
\end{notesbox}

\begin{extensionbox}
\textbf{Finding eigenvalues with a determinant (a first look at Lesson~6.1).} An eigenvector satisfies
$(A-\lambda I)\xx=\zero$ for a nonzero $\xx$, so $A-\lambda I$ must be singular: $\det(A-\lambda I)=0$.
Use $A=\begin{bmatrix} 3 & 1 \\ 1 & 3 \end{bmatrix}$.
\begin{enumerate}[label=(\alph*), leftmargin=1.8em, itemsep=5pt]
  \item Write $A-\lambda I$ and compute $\det(A-\lambda I)$ as an expression in $\lambda$.
  \item Set it to $0$ and solve for $\lambda$. (You should get two values.)
\end{enumerate}
\writelines{3}
\end{extensionbox}

\begin{spiralbox}
\textbf{Coming next --- Lesson 6.1: Introduction to Eigenvalues.} Today you \emph{tested} vectors you
were handed. Next we \emph{find} the special directions from scratch: solve $\det(A-\lambda I)=0$ for
the eigenvalues $\lambda$, then solve $(A-\lambda I)\xx=\zero$ for each eigenvector --- no guessing.
\end{spiralbox}

\end{document}
